% Figure from MATH 551 notes, section 18. Compile with the notes preamble.
\begin{tikzpicture}
\begin{axis}[nfig, width=0.70\textwidth, height=0.50\textwidth, clip=false,
  axis lines=middle, xmin=-0.42, xmax=0.46, ymin=-0.42, ymax=0.44,
  xtick=\empty, ytick=\empty, xlabel={}]
  \addplot[black!50, dashed, thin, domain=0:0.4] {x};
  \addplot[black!50, dashed, thin, domain=0:0.4] {-x};
  \addplot[figB, thick, domain=-0.4:0, samples=2] {0};
  \addplot[figB, thick, domain=0.004:0.4, samples=3000] {x*sin(deg(1/x))};
  % one right secant
  \addplot[figR, thick] coordinates {(0,0) (0.2457,-0.1966)};
  \addplot[only marks, mark=*, mark size=1.3pt, figR] coordinates {(0.2457,-0.1966)};
  \addplot[only marks, mark=*, mark size=1.4pt, black] coordinates {(0,0)};
  \node[font=\scriptsize, black!65, anchor=west] at (axis cs:0.40,0.40) {slope $1 = D^+f(0)$};
  \node[font=\scriptsize, black!65, anchor=west] at (axis cs:0.40,-0.40) {slope $-1 = D_+f(0)$};
  \node[font=\scriptsize, figB, anchor=south] at (axis cs:-0.21,0.02) {$D^-f(0) = D_-f(0) = 0$};
  \node[font=\scriptsize, figR, anchor=north west] at (axis cs:0.3,-0.2855) {$(h, f(h))$};
  \node[font=\small, figB, anchor=south] at (axis cs:0.37,0.25) {$f$};
\end{axis}
\end{tikzpicture}
